\section{Annexe B --- Anatomie d'un trade}
\label{app:signal_flow}

\lettrine[lines=2, findent=2pt, nindent=0pt]{P}{o}ur ancrer la compréhension opérationnelle de la stratégie, cette annexe déroule pas à pas le cheminement qui mène à une entrée réelle sur le Moteur~1 (la composante principale du portefeuille). Les mécaniques des Moteurs~2 et~3 suivent la même logique de filtrage mais sur horizon journalier~; elles sont détaillées dans les sections~\ref{sec:ts_momentum} et~\ref{sec:rsi_daily}.

\subsection*{Exemple : entrée long sur EUR-USD le 15~juin 2024 à 10h30~UTC}

\begin{enumerate}
\item \textbf{Calcul du VWAP ancré au jour courant} depuis l'ouverture à 00h00~UTC jusqu'à la minute courante.
\item \textbf{Calcul de la déviation} entre le prix de clôture actuel et le VWAP, puis comparaison à la bande inférieure calculée à partir de l'écart-type glissant sur les 80 dernières minutes.
\item Si la déviation est \emph{sous} la bande inférieure, cela signale une sur-vente intraday exploitable.
\item \textbf{Filtre session}~: l'heure courante 10h30~UTC est bien dans la fenêtre de chevauchement Londres~+~New~York (6h--14h~UTC) où la liquidité est maximale.
\item \textbf{Filtre macro premier étage}~: le spread 10Y--2Y des bons du Trésor américain à cette date est inférieur à 0.5, donc pas de stress cyclique majeur.
\item \textbf{Filtre macro second étage}~: le taux de chômage américain sur trois mois n'est pas en hausse, donc pas de fin de cycle d'expansion.
\item Les trois filtres sont passants, et la position longue est ouverte.
\item La taille de position est calculée à partir du pool de trésorerie alloué au Moteur~1 (80\,\% du portefeuille total) multiplié par le levier global courant, divisé par le prix d'entrée.
\end{enumerate}

\textbf{Conditions de sortie (la première atteinte déclenche la clôture)~:}
\begin{itemize}
\item Retour du prix au-dessus du VWAP (take-profit souple)~;
\item Ou stop-loss à $-0.5\,\%$ du prix d'entrée~;
\item Ou take-profit dur à $+0.6\,\%$ du prix d'entrée~;
\item Ou durée maximale de détention de 6~heures dépassée~;
\item Ou arrêt forcé à 21h00~UTC (fin de session).
\end{itemize}

\vfill

\begin{center}
\color{primary}\sffamily\Large\textbf{--- Fin du rapport ---}\\[0.5em]
{\color{slateGrey}\normalsize\itshape
Pour toute question méthodologique, demande d'extension ou analyse complémentaire,\\
le lecteur est invité à contacter directement l'auteur.}
\end{center}
